\documentclass[11pt]{article}

\usepackage[margin=1in]{geometry}

\usepackage{booktabs}

\usepackage{graphicx}

\usepackage{float}

\usepackage{adjustbox}

\usepackage{caption}

\captionsetup{font=small,labelfont=bf,skip=4pt}

\title{LLM-as-a-Judge Report: Repair-understanding\\\large GSM8K}

\date{}

\begin{document}

\maketitle

\section*{What is measured?} Repair-understanding is the judge's assessment of how much the model loses the intended meaning of the question because of the corrupted wording. A score of 0 means the meaning is fully understood; 5 means the model is fundamentally lost. This is not a label for whether the arithmetic answer is correct.

\section*{How to read the score} The judge uses a 0--5 scale. The reported value in the tables and graphs is the arithmetic mean over valid judged traces in the group. These are behavioral scores, not correctness scores: a model can be wrong without showing doubt, and it can be correct while expressing uncertainty.

\section*{Examples} The typo-percentage trend supports the requested direction: the pooled mean increases from 0.053 at 0\% to 0.580 at 25\%, 0.576 at 50\%, and 0.764 at 75\%. The real-word ratio comparison is less monotonic: it is 0.554 at 10\%, 0.697 at 40\%, and 0.669 at 70\%. Thus more typo corruption is associated with lower understanding in this sample, while the real-word-ratio effect peaks at 40\% rather than increasing through 70\%.

\begin{table}[H]
\centering
\caption{Examples from the per-configuration aggregate table. The relevant score is shown alongside both judge dimensions for context.}
\small
\begin{tabular}{lrrrrr}
\toprule
\textbf{config} & \textbf{n} & \textbf{selected score} & \textbf{self-doubt} & \textbf{repair-understanding} \\
\midrule
clean & 494 & 0.053 & 0.053 & 0.053 \\
typo75\_real10 & 482 & 0.660 & 0.660 & 0.660 \\
typo75\_real70 & 486 & 0.817 & 0.817 & 0.817 \\
\bottomrule
\end{tabular}
\end{table}

\section*{Aggregate result by typo percentage} \begin{table}[H]
\centering
\caption{Mean score pooled by typo percentage. Clean is the 0\% baseline; typo points pool all three real-word ratios.}
\small
\begin{tabular}{lrrr}
\toprule
\textbf{typo \%} & \textbf{n} & \textbf{mean score} \\
\midrule
0 & 494 & 0.053 \\
25 & 1450 & 0.580 \\
50 & 1454 & 0.576 \\
75 & 1455 & 0.764 \\
\bottomrule
\end{tabular}
\end{table}

\begin{figure}[H]\centering\includegraphics[width=0.82\textwidth]{graphs/repair\_understanding\_score\_by\_typo\_rate.png}\caption{Repair-understanding pooled by typo percentage.}\end{figure}

\section*{Aggregate result by real-word ratio} \begin{table}[H]
\centering
\caption{Mean score pooled by real-word ratio. Each point pools the 25\%, 50\%, and 75\% typo configurations.}
\small
\begin{tabular}{lrrr}
\toprule
\textbf{real-word ratio \%} & \textbf{n} & \textbf{mean score} \\
\midrule
10 & 1447 & 0.554 \\
40 & 1457 & 0.697 \\
70 & 1455 & 0.669 \\
\bottomrule
\end{tabular}
\end{table}

\begin{figure}[H]\centering\includegraphics[width=0.82\textwidth]{graphs/repair\_understanding\_score\_by\_real\_ratio.png}\caption{Repair-understanding pooled by real-word ratio.}\end{figure}

\section*{Conclusion}

The typo-percentage trend supports the requested direction: the pooled mean increases from 0.053 at 0\% to 0.580 at 25\%, 0.576 at 50\%, and 0.764 at 75\%. The real-word ratio comparison is less monotonic: it is 0.554 at 10\%, 0.697 at 40\%, and 0.669 at 70\%. Thus more typo corruption is associated with lower understanding in this sample, while the real-word-ratio effect peaks at 40\% rather than increasing through 70\%.

\end{document}